Table [] reports the estimates by school types. Column one reports estimates for the preferred specification of a Tobit model. I find independent schools assigned on average 0.11 standard deviation points higher grades to their students than academy schools. Further education schools assigned 0.72 standard deviation points lower grades than academy schools. Column 2 reports the estimates using an OLS specification. The estimates for independent schools still remain positive, although the magnitude is around half of the Tobit specification. I infer this is due to the censoring of the dependent variable. Since many independent school are well equipped with educational material, an OLS estimates miss-measures the impact of the treatment by under-estimating the effect on students in affluent school. Column 3 reports the coefficient estimate for the probit model using A or A$^*$ as the cutoff. The signs of the coefficients follow the coefficients of the Tobit model. The estimate of sixth form college turns negative because their student's grades did not reach A or $A^*$ as much as academy schools. 

I find that independent schools inflated their students more than other type of schools. I infer this is due to application patterns and the source of funding of the independent schools. Independent schools run by collecting tuition from student's parents, while other types of schools collect funding from the state. Since parents often choose schools by the schools placement records, independent schools aim to improve their placement records. Also students from independent schools apply to reach schools more often, therefore inflating the grades of students increases the chances of their students landing in elite universities. Students in public schools in the other hand, applies to less prestigious universities and teachers are not as incentived to place them in elite institutions. Despite that grammar schools are selective with their students, I do not find any effect of them giving grades more generously than academy schools. I find further education schools did not inflate their students grade as much as others. 

Table [] reports relation between the degree of inflation and school quality. The first column reports the estimates from the Tobit model. I find a non-linear reverse U shaped relation between school quality and the degree of inflation. Strong and Strong$^*$ denote the share of students with strong GCSE pass scores. I find that schools with higher performing students inflated less that those with lower performing students. The decreasing rate is convex, indicating that schools in the middle range of the quality distribution inflated students more than others\footnote{Show distribution of the school quality}. Schools at the bottom end of the school quality distribution also inflated. \emph{Normal} denotes the share of student that passed GCSE with moderate grades. I find that worse schools inflated more than better schools. Column 2 reports the estimates using an OLS specification. The OLS estimates the schools with more students with high GCSE scores did not inflate their scores at all. This is due to the censoring of the dependent variable, since the students in good secondary schools cannot be inflated because their scores gets under measured due to the truncated distribution of grades. The results for the \emph{Normal} measure also highlights how the degree of inflation is shrinked towards zero by using OLS. Column 3 reports the Probit model using a A/$A^*$ cutoff as the dependent variable. After adjusting the estimates with marginal effects, which I condition the estimates at zero covariate values, the results indicate that better schools were more likely to receive an A/$A^*$ (Figure []). However, the school quality measure using the \emph{Normal} measurements predicts opposite trends for schools at the top end of the average score distribution. As each measurement of school distribution is strongly positively correlated at the top of the distribution, the marginal effect calculation for the top schools may require using the two measurements together for correctly estimating the marginal effect. The estimates indicate better schools were more likely to receive A/$A^*$ than worse schools, even after controlling for students past GCSE performance.

I infer from the  that the distribution of grades within a classroom restricts the degree of inflation that is possible for teachers. I hypothesize that teachers were morally constraint on preserving the rank of their students, even after the grading method switch. Teachers may have adverse preference over re-ranking students even those the grading scheme did not constraining them to inflate all their students grades. If that were the case, then low performing students in high-performing schools with many high performing students cannot benefit from grade inflation as much as a similarly performing student in a low performing school. Low performing schools can inflate a majority of their students without changing the rank of their students, while in high performing schools teachers are upper bounded by the extent that they can inflate low performing students grades because that would lead to reversing the pre-existing rank of students. 

The results from the binary dependent variable indicates that although the grade inflation was more intense in worse schools, the increments in the number of top grades were still concentrated in the top schools. I attribute this to differences across schools with respect to students score distance to the A/$A^*$ threshold. Even accounting for GCSE score, schools differ in terms of equipments and teacher quality that positively contributes to the students probability of receiving an A/$A^*$ grade at their A-levels. This underscores that although the grade inflation quantitatively benefited lower perfoming schools than higher one, the population of students receiving A/$A^*$, which arguably is the grade for securing a placement at elite schools, was still concentrated in good schools.

Table [] decomposes the degree of inflation by school type and their student quality. I find independent schools and state schools inflated grades regardless of their pupil's grades. Independent schools assigned 0.15 standard division higher scores, and state schools assigned 0.09 standard deviation points higher scores to  their pupils than academy schools. 

Table [] reports the difference in degrees of inflation between subject groups. I find a 10 percentage point difference between facilitating and non-facilitating subjects with respect to their probability of receiving an A or A$^*$. However, the estimate for facilitating subjects is positive at 6 percentage points. Column (2) reports the estimate for quantitative subjects against non-quantitative subjects. I find that the the gap between the two categories is larger than the facilitating subject categorization. Non-quantitative subjects were 13 percentage points more likely to receive an A or A$^*$ than non-quantitative subjects. Non-quantitative subjects improved their probability of receiving a top grade by 5 percentage points, which is not large as the facilitating subjects. I attribute the differences across subject types regarding their degree of inflation to teacher's discretion over grading. Teaching quantitative subjects often subjects students to take internal exams with readily marked questions. Teachers learn about their students capability with future A-levels accurately. On the other hand, subjective subjects, such as drama studies or media studies, burdens teachers to rank student based on performance outcomes with opaque quality. The different degrees of subjectiveness across subject groups contributed to differential impact of the grading method change.

The results indicate student's grades were differentially shifted upwards based on their subject choice, possibly facilitating students from marginalized population to enter their reach schools. Most quantitative subjects are included in the category of \emph{facilitating subject}, a group of subjects that elite universities recommend their applicants to include in their applications. Independent schools offer more competitive courses to their students to increase their student's chances of acceptance, and also attract more students. The differential shift in grade standards across subjects likely provided more equity to students in less affluent backgrounds. 

Table [] reports the relation between student demographics of schools and the degree of inflation. I find a positive relation between the female share of the school and the degree of inflation. A one percentage point increase in the female share is associated with 0.0007 standard deviation point increase in the grade assigned. I find a negative relationship between non-while share of students and the extent of inflation. A one percentage point increase in the share of white students is associated with 0.002 standard deviation point increase in the grade assigned. I find a positive relationship between foreign student share and the degree of inflation, however the point estimate is unstable. I find a positive relationship between share of disabled students and degree of inflation. I find a positive relationship between average number of A-level related exams taken by students. Schools taking more A-level or AS-level subjects received higher grades.  A one percentage point increase in the average number of subjects is associated with 0.02-0.04 standard deviation point increase in the grade assigned.

Figure [] presents the extent of grade inflation across different score bins in their previous GCSE. All bin groups experience a 0.4 standard deviation increase in their grades. The results hold for different subjects of the GCSE.

Table [] reports the acceptance and rejection rate changes. I find the number of rejection per a student decreased by 0.02 application due to the grade inflation. The decrease is significantly larger for independent schools. Independent schools decreased the number of rejection the most out of all type of schools, although the sign of the difference is not statistically significant. Further education student did not change the number of rejections.

Table [] reports the number of acceptance by university class. Students were accepted into a higher tier university with higher probability than their previous cohorts. The beneficiaries of the were students from academy schools, Further education schools, sixth form colleges, and state schools. All 4 types of schools experienced a positive increase in their probability of getting accepted. Grammar schools and independent schools did not change their numbers. However, independent schools did improve the number of students getting accepted in mid-tier tariff schools. Whereas for other schools it was a net decrease. 

Despite the grade method change shifted the grade distribution upwards across a diverse body of students, students that increased the rank of accepted schools through the loosee grade standards were limited, by the beneficiaries concentrated in schools associated with affluent backgrounds such as independent schools. I attribute this to application pattern of students which was set before the pandemic happened. Firstly, students from independent schools are more to apply to better universities, that is out of their possible range, possibly due to their peers applying to better universities and teachers urging them to apply higher. Secondly, students from marginalized population in the British higher education system may not prefer to apply to higher ranked universities due to concerns on immersing with their peers at university. 

The results also underscore how swiftly schools and students adjust to changes in grading standards. Despite the significant change in the grades, university programs could not react by raising their admission standards in similar degrees.